%=======================================================================
% Six-Yao Sensor Arrays: A Signal-Agnostic Yijing-Inspired Framework
% for Unsupervised Industrial Image Anomaly Detection
% Draft v0.1 (2026-09-13)
% NOTE: numbers marked [TBD-15class] pending full MVTec AD (15 classes)
%=======================================================================
\documentclass[journal,10pt]{IEEEtran}
\usepackage{amsmath,amssymb,amsfonts}
\usepackage{graphicx}
\usepackage{booktabs}
\usepackage{multirow}
\usepackage{cite}
\usepackage{algorithm}
\usepackage{algpseudocode}
\usepackage{url}
\usepackage{xcolor}

\newcommand{\tbd}[1]{\textcolor{red}{[TBD: #1]}}

\title{A Yijing Yao-Variation Framework for Unsupervised Industrial Image
Anomaly Detection}

\author{%
\IEEEauthorblockN{Xinglu Ma, \textit{Member, IEEE}, et al.}
\IEEEauthorblockA{School of Information Science and Technology\\
Qingdao University of Science and Technology, Qingdao, China\\
Email: \texttt{mxl@qust.edu.cn}}
}

\begin{document}
\maketitle

\begin{abstract}
Industrial image anomaly detection (AD) is dominated by methods that
rely on ImageNet-pretrained features and large corpora of normal samples,
which hinders deployment on edge devices and obscures interpretability.
Departing from the \emph{Yijing} (易经), we propose an \emph{unsupervised}
industrial image AD framework whose theoretical core is the idea of
\emph{yao variation} (爻变): the essence of an anomaly is not a deviation
in absolute magnitude, but a change in \emph{relational structure}. We
therefore map the anomaly-response field of an image---or of its deep
features---to the readings of a \emph{six-yao sensor array}, and declare an
anomaly when the \emph{correlation structure} among the yao readings
departs from its normal manifold, rather than when a geometric
(orthogonal) distance grows. The framework is \emph{signal-agnostic}: the
same Yijing criterion is instantiated with (i) purely hand-crafted signals
requiring \emph{no pretraining}, and (ii) frozen ImageNet-pretrained
ResNet features. On all 15 classes of MVTec AD, the pretrained
instantiation attains an average image-level AUROC of \textbf{0.862}
(deep-layer six-yao sensors), approaching the reported PaDiM reference
(\textbf{0.975}) and PatchCore (\textbf{0.992})~\cite{defard2021padim,roth2022patchcore},
while the hand-crafted instantiation reaches \textbf{0.709} with \emph{no}
pretraining at all and full interpretability.
A controlled ablation over feature extractors on a common 4-class subset
(\texttt{bottle}, \texttt{tile}, \texttt{metal\_nut}, \texttt{toothbrush}):
random-initialized ResNet \textbf{0.772} $<$ hand-crafted fixed convolution
\textbf{0.799} $<$ hand-crafted six-yao signals \textbf{0.854} $<$
ImageNet-pretrained ResNet \textbf{0.932}.
On the full 15 classes the hand-crafted-six-yao and pretrained
instantiations reach \textbf{0.709} and \textbf{0.862}. This shows that the framework does
not \emph{produce} knowledge but \emph{organizes} it: performance is bounded
by the information content of the signal. The framework learns online from
the very first sample via Welford updates with an auto-adaptive threshold,
is robust to cold-start contamination, and is data-efficient.
\end{abstract}

\begin{IEEEkeywords}
Industrial anomaly detection, unsupervised learning, Yijing, yao
variation, six-yao sensor array, interpretability, online learning,
MVTec AD.
\end{IEEEkeywords}

%=======================================================================
\section{Introduction}
\IEEEPARstart{I}{ndustrial} visual anomaly detection aims to identify
defective items in a production line using only normal samples for
modeling, since defects are rare, diverse, and unpredictable. The
MVTec AD benchmark~\cite{bergmann2019mvtec} has become the de facto
standard, and recent state-of-the-art (SOTA) methods---e.g., PaDiM
~\cite{defard2021padim} and PatchCore~\cite{roth2022patchcore}---report
image-level AUROC above 0.95 across its 15 categories.

Despite their strong performance, these methods share two limitations.
First, they \emph{depend on ImageNet-pretrained features}; the visual
prior encoded in such weights is not always available or desirable, and
its contribution is rarely isolated in a clean ablation. Second, they
organize features via statistical distances or memory banks
($k$-nearest-neighbor lookups), which are difficult to interpret and
require storing large reference sets at inference time.

In this work we take a different organizing principle, inspired by the
\emph{Yijing} (易经, \emph{Book of Changes}). The \emph{Yijing} models the
state of the world through six \emph{yao} (爻, lines), and---crucially---
characterizes change through \emph{yao variation} (爻变): the transition
of yao and the alteration of their mutual relations. We translate this
into an anomaly-detection framework whose theoretical core is captured in
one sentence: \emph{an anomaly lies not in a deviation of magnitude, but
in a change of relations}. Accordingly, the deviation field of an image is
read by a \emph{six-yao sensor array}, and an anomaly is declared when the
\emph{correlation structure} among the sensors' readings departs from its
normal manifold. We deliberately avoid orthogonality-based criteria:
correlation among yao is treated as \emph{information} to be exploited,
not redundancy to be removed.

Our framework is \emph{signal-agnostic}: the six-yao operators act on a
residual field $R = |F - T| / \mathrm{SD}$ built from \emph{any} feature
representation $F$ and an online template $T$. We instantiate it on two
signal sources: hand-crafted image statistics (no pretraining) and frozen
ImageNet-pretrained ResNet features. The same online engine---trained
incrementally from the first sample with an auto-adaptive threshold---
serves both.

\subsection{Contributions}
\begin{itemize}
\item \textbf{Six-yao sensor model.} We map any anomaly-response field to
six interpretable readings, decoupling the detection framework from the
feature extractor.
\item \textbf{Yao-variation criterion.} We judge anomalies by changes in
inter-yao correlation structure and explicitly reject orthogonalization;
we report this as a controlled (negative) result.
\item \textbf{Online, parameter-free adaptation.} Welford incremental
statistics, an auto-adaptive threshold (the $q{=}75$ quantile of normal
scores), and robustness to cold-start contamination.
\item \textbf{Signal-agnosticism and a knowledge-vs-carrier ablation.}
On the same framework, hand-crafted signals reach \textbf{0.709} and
pretrained features \textbf{0.862} on all 15 MVTec AD classes; a controlled
sweep over extractors (random $<$ fixed-conv $<$ pretrained) shows
performance is bounded by signal information, not by the framework.
\item \textbf{Data efficiency.} Strong few-shot behavior; the pretrained
instantiation recovers most of its full-data AUROC from only 10 normal
images per class.
\item \textbf{Edge-friendliness.} No memory bank: only a template is
stored, giving \textbf{$52\times$} less storage and \textbf{$55\times$} faster
inference than a PatchCore-style $k$NN bank.
\end{itemize}

%=======================================================================
\section{Related Work}
\subsection{Industrial Anomaly Detection}
MVTec AD~\cite{bergmann2019mvtec} established a 15-category benchmark for
unsupervised AD. \textbf{PaDiM}~\cite{defard2021padim} models patch-wise
feature distributions with multivariate Gaussians and Mahalanobis
distance. \textbf{PatchCore}~\cite{roth2022patchcore} uses a greedy coreset
of pretrained embeddings and nearest-neighbor scoring. \textbf{SPADE}
~\cite{cohen2021spade} retrieves on a normal-image feature bank.
More recent lines include zero-/few-shot detection with vision-language
models, e.g.\ \textbf{WinCLIP}~\cite{jeong2023winclip}, and
synthetic-anomaly based methods, e.g.\ \textbf{SimpleNet}
~\cite{liu2023simplenet} and \textbf{DRAEM}~\cite{zavrtanik2021draem};
these improve accuracy but generally increase model complexity and
storage. In contrast, our framework is training-free for the pretrained
path, stores only a template, and is architecture-agnostic.

\subsection{Hand-crafted and Traditional Methods}
Before deep features became dominant, industrial inspection relied on
hand-crafted descriptors---Gabor/filter banks, Local Binary Patterns
(LBP), gray-level co-occurrence matrices (GLCM), and structural-texture
features---combined with statistical classifiers. These methods are
interpretable and cheap but struggle with the diversity of real defects.
Intensity- and frequency-domain residuals against a reference template
remain common in classical surface inspection. Our hand-crafted
instantiation is in this lineage, but differs by (i) organizing responses
into a fixed, interpretable six-yao structure and (ii) learning the normal
statistics online with an auto-adaptive threshold.

\subsection{Online / Incremental Unsupervised Learning}
Incremental statistics (Welford) and concept-drift adaptation are
classical topics in streaming data analysis. In anomaly detection,
few-shot and continual-learning settings have recently received
attention~\cite{jeong2023winclip}. Unlike memory-bank methods that fix a reference
set after a training phase, our engine updates its template, scale, and
yao statistics incrementally and is therefore naturally robust to
cold-start contamination.

\subsection{Yijing and Computing}
The \emph{Yijing} (\emph{Book of Changes}) has inspired prior
cross-disciplinary work connecting its hexagram/yao structure with
computation, coding, and machine learning~\cite{zhang2020yijing}. To our knowledge,
this is the first work to operationalize the \emph{yao-variation}
principle as a concrete, competitive unsupervised anomaly-detection
criterion.

%=======================================================================
\section{Theoretical Framework}
\label{sec:framework}
\subsection{The Six-Yao Sensor Model}
Let $F$ denote a feature field of an image and $T$ an online template.
We define the \emph{anomaly-response field}
\begin{equation}
R = \frac{|F - T|}{\mathrm{SD}},
\end{equation}
where $\mathrm{SD}$ is the per-location standard deviation of $F$ over
normal samples. A \emph{six-yao sensor array} is a map
$\mathcal{S}: R \mapsto \mathbb{R}^{H\times W\times 6}$, assigning to each
spatial location six readings $s_1,\dots,s_6$ (初爻, 二爻, \dots, 上爻),
each capturing a distinct mode of anomaly response (Table~\ref{tab:yao}).

\begin{table}[t]
\centering
\caption{Six-yao operator definitions.}
\label{tab:yao}
\begin{tabular}{cl}
\toprule
Yao & Operator (on $R$) \\
\midrule
初 (1) & residual peak: $\max(R)$ \\
二 (2) & high-frequency FFT ratio \\
三 (3) & structural-tensor anisotropy (99th pct.) \\
四 (4) & $\max|R-\bar R| / \mathrm{std}(R)$ \\
五 (5) & color residual peak \\
六 (6) & high-contrast connected-region density \\
\bottomrule
\end{tabular}
\end{table}

\subsection{Yao Variation as the Anomaly Criterion}
\label{sec:yaobian}
The central departure from conventional statistical AD is the criterion
itself. Let $\mathbf{s} \in \mathbb{R}^6$ be the six readings at a
location, with online normal statistics $\mu, \sigma$ and online
correlation matrix $C$. We define the standardized reading
\begin{equation}
z_i = \left|\frac{s_i - \mu_i}{\sigma_i}\right|,
\end{equation}
and score a location by combining \emph{marginal deviation} with the
\emph{variation of inter-yao relations}:
\begin{equation}
\mathrm{score} = \max_i z_i \quad\text{modulated by}\quad
\Delta C = \|C(R) - C_{\text{normal}}\|.
\end{equation}
We emphasize that the six yao are \emph{not} orthogonalized: correlation
among yao readings carries discriminative information. In
Sec.~\ref{sec:ablation} we show empirically that PCA-whitening the yao
space \emph{degrades} AUROC, confirming that the yao-variation view is the
correct one.

\subsection{Auto-Adaptive Threshold}
Rather than fixing a decision threshold manually, we set it to the
$q$-th quantile of the model's own normal scores, with $q{=}75$
(a parameter-free, robust operating point allowing one false alarm per
four normal items). We show in Sec.~\ref{sec:ablation} that no single
quantile rule is optimal across all categories, motivating this fixed,
principled choice.

\subsection{Online Learning}
All task-specific statistics---template $T$, scale $\mathrm{SD}$, and the
yao statistics $\mu,\sigma,C$---are updated incrementally via Welford's
algorithm from the very first (unlabeled) sample:
\begin{equation}
\bar{x}_n = \bar{x}_{n-1} + \frac{x_n - \bar{x}_{n-1}}{n}, \quad
M_n = M_{n-1} + (x_n - \bar{x}_{n-1})(x_n - \bar{x}_n).
\end{equation}
The $1/n$ weighting induces a natural forgetting that makes the model
robust to cold-start contamination (Sec.~\ref{sec:ablation}).

%=======================================================================
\section{Method}
\subsection{Overall Pipeline}
\begin{figure}[t]
\centering
\includegraphics[width=\columnwidth]{figures_paper/fig1_framework.png}
\caption{Overview of the six-yao sensor framework. An
input image yields a residual field $R=|F-T|/\mathrm{SD}$ from any feature extractor $F$ and an
online template $T$; six yao sensors read $R$ into a six-dimensional
reading; an anomaly is declared when the inter-yao correlation structure
departs from its normal manifold.}
\label{fig:framework}
\end{figure}
The overall pipeline is shown in Fig.~\ref{fig:framework}.
Given an image, we compute a feature field $F$, build the residual field
$R = |F-T|/\mathrm{SD}$ against the online template, extract six-yao
readings $\mathcal{S}(R)$, standardize them with online statistics, and
take the image-level score as the maximum over locations and yao. An
anomaly is declared if the score exceeds the auto-adaptive threshold.

\subsection{Six-Yao Operators}
\begin{figure}[t]
\centering
\includegraphics[width=\columnwidth]{figures_paper/fig2_yao_operators.png}
\caption{The six yao operators acting on the residual field $R$:
(1) local peak (top-hat), (2) high-frequency energy,
(3) residual gradient, (4) channel dispersion, (5) channel energy, (6) local
aggregation.}
\label{fig:yao}
\end{figure}
We instantiate the six operators (Fig.~\ref{fig:yao}) as in
Table~\ref{tab:yao}. Operator~(1)
captures localized high-contrast residuals; (2) captures texture/frequency
irregularities; (3) captures directional structure breaks; (4) captures
non-Gaussian residual structure; (5) captures chromatic anomalies; and
(6) captures connected high-contrast regions. We adopt \emph{max-type}
statistics (e.g., $\max R$, $\max|R-\bar R|/\mathrm{std}$) rather than
percentiles, as they better capture localized defects.

Concretely, let $R_c=\max_c R$ and $R_m=\mathrm{mean}_c R$ be the channel-wise
maximum and mean of the residual field. With a local background estimate
$b=\mathrm{uniform}_7(R_c)$ (a $7{\times}7$ uniform filter), the six readings
are
\begin{align}
s_1 &= \max(R_c-b,\,0), && \text{(top-hat: local peak)}\\
s_2 &= |\nabla^2 s_1|, && \text{(high-frequency)}\notag\\
s_3 &= \sqrt{(\partial_y R_c)^2+(\partial_x R_c)^2}, && \text{(gradient)}\notag\\
s_4 &= \mathrm{std}_c(R), && \text{(channel dispersion)}\notag\\
s_5 &= s_1^2, && \text{(local energy)}\notag\\
s_6 &= \sqrt{\max(b_2-b_1^2,0)},\ b_2{=}\mathrm{unif}_5(R_c^2),\ b_1{=}\mathrm{unif}_5(R_c). \notag
\end{align}
The top-hat in $s_1$ highlights locally salient but absolutely small peaks
(fine scratches), and $s_6$ measures local concentration via a windowed
variance.

\subsection{Signal Instance 1: Hand-Crafted (Zero Pretraining)}
In this instance the feature field is the image itself: $F$ is the resized
RGB image (and its grayscale), and the template $T$ is an online mean
image. The residual field $R=|F-T|/\mathrm{SD}$ therefore measures
deviation from the running normal appearance. The six-yao operators are
instantiated directly on $R$ (denoting $R_b$ as a Gaussian-smoothed
residual):
\begin{align}
s_1 &= \max R_b, && \text{(peak / \emph{ji})}\\
s_2 &= \frac{\sum_{\text{high}}|\hat R|}{\sum|\hat R|}, &&
        \text{(high-freq.~{}ratio)}\\
s_3 &= \mathrm{P}_{99}\!\left(\frac{\sqrt{(J_{xx}{-}J_{yy})^2+4J_{xy}^2}}{J_{xx}+J_{yy}}\right), &&
        \text{(structure-tensor anisotropy)}\notag\\
s_4 &= \frac{\max|R_b-\bar R_b|}{\mathrm{std}(R_b)}, && \text{(kurtotic deviation)}\notag\\
s_5 &= \max R_c, && \text{(color residual peak)}\notag\\
s_6 &= N\,(1+\bar\kappa)\,\max R_b, && \text{(high-contrast region density)}\notag
\end{align}
where $\hat R=\mathcal{F}(R_b-\bar R_b)$ is the Fourier transform, $J$ is
the structure tensor (with Gaussian-windowed products $J_{uv}=\mathcal{G}_2(\partial_u R_b\,\partial_v R_b)$), $\mathrm{P}_{99}$ is the 99th
percentile, $R_c$ is the color residual, $N$ is the number of connected
components above a contrast threshold, and $\bar\kappa$ is their mean
normalized perimeter (compactness). This instance uses \emph{no}
pretrained weights and is fully interpretable; it also matches the
``hand-crafted six-yao signals'' rung of the ladder in
Sec.~\ref{sec:ablation}.

\subsection{Signal Instance 2: Pretrained ResNet Features}
We use ResNet18~\cite{he2016resnet} with ImageNet-1k weights
(IMAGENET1K\_V1) as a \emph{frozen} feature extractor (no fine-tuning, no
target-domain labels). We take layer1 ($/4$, 64ch), layer2 ($/8$, 128ch),
and layer3 ($/16$, 256ch), upsample to a common $/4$ grid, and apply the
six-yao operators to each layer. We stress that the extractor is frozen
and the adaptation is purely online; the method is \emph{not} zero-shot
(ImageNet weights are a prior) but requires no target-domain supervision.

\subsection{Multi-Layer Sensor Fusion}
Layer1 ($/4$) captures fine local structure, while layer3 ($/16$)
captures shape/semantic anomalies. Averaging their per-layer six-yao
scores (``L1+L3'') consistently outperforms any single layer; we use this
as the default pretrained configuration.

\subsection{Online Engine}
Algorithm~\ref{alg:engine} summarizes the online engine. The template and
scale are updated in a single pass over the (unlabeled) stream; the yao
statistics $\mu,\sigma$ are maintained by Welford updates; the decision
threshold is recalibrated as the $q{=}75$ quantile of the current normal
scores (re-scoring a small cached normal subset with the mature model,
which avoids the immaturity of early estimates).

\begin{algorithm}[t]
\caption{Online six-yao anomaly engine}
\label{alg:engine}
\begin{algorithmic}[1]
\State \textbf{init} $T\!\leftarrow\!0$, $\mathrm{SD}\!\leftarrow\!1$,
       $n\!\leftarrow\!0$
\For{each incoming sample $x$}
    \State $F\leftarrow\mathrm{features}(x)$
    \If{$n<2$} $T\leftarrow F$, $\mathrm{SD}\leftarrow 1$, $n\!+\!+$
    \Else
        \State $n\!+\!+$;
               $T\leftarrow T+\frac{F-T}{n}$;
               $\mathrm{SD}\leftarrow\sqrt{\smash{(1-\tfrac1n)\mathrm{SD}^2+\tfrac1n(F-T)^2}}$
    \EndIf
    \State $R\leftarrow|F-T|/\mathrm{SD}$;
           $\mathbf{s}\leftarrow\mathcal{S}(R)$
    \State $\mu\leftarrow\mu+\frac{\mathbf{s}-\mu}{n}$;
           $\sigma\leftarrow\mathrm{update}$ \Comment{Welford}
    \State $z\leftarrow|\mathbf{s}-\mu|/\sigma$
    \State $\mathrm{score}(x)\leftarrow\max_{p,i} z[p,i]$
    \State $\tau\leftarrow\mathrm{quantile}_{75}\{\mathrm{score}\text{ of normal subset}\}$
    \State \textbf{return} \texttt{defect} if $\mathrm{score}(x)>\tau$ else \texttt{normal}
\EndFor
\end{algorithmic}
\end{algorithm}

\subsection{Complexity and Edge-Friendliness}
The engine stores only a template and its scale (a single feature-map-sized
array) plus a $6{\times}6$ yao covariance---\emph{O}(1) in the number of
normal samples---whereas memory-bank methods such as PatchCore retain a
large coreset of embeddings and perform nearest-neighbor search at
inference. We quantify this on MVTec \texttt{bottle} (CPU): our sensor
array requires \textbf{29.4\,MB} of stored reference data and
\textbf{66\,ms} per image, versus \textbf{1534\,MB} and
\textbf{3678\,ms} for a PatchCore-style patch-$k$NN bank ($856{,}064$
normal patches)---a \textbf{$52\times$} storage and \textbf{$55\times$}
inference-speed advantage (Table~\ref{tab:cost}, Fig.~\ref{fig:light}).
Inference is a single forward pass plus element-wise yao operators; no
$k$NN lookup or large matrix inversion is required.

\begin{table}[t]
\centering
\caption{Lightweightness comparison (MVTec \texttt{bottle}, CPU).}
\label{tab:cost}
\begin{tabular}{lccc}
\toprule
Method & Storage (MB) & Inference (ms) & Reference set \\
\midrule
Six-yao array (ours) & \textbf{29.4} & \textbf{66} & template only \\
patch-$k$NN (PatchCore) & 1534 & 3678 & 856{,}064 patches \\
PaDiM-style (diag.) & 51.1 & 43 & per-channel stats \\
\bottomrule
\end{tabular}
\end{table}

\begin{figure}[t]
\centering
\includegraphics[width=\columnwidth]{figures_paper/fig12_lightweight.png}
\caption{Lightweightness: storage and inference latency of the six-yao
array versus a PatchCore-style patch-$k$NN bank (MVTec \texttt{bottle},
CPU). Our method stores $52\times$ less and runs $55\times$ faster.}
\label{fig:light}
\end{figure}

%=======================================================================
\section{Experiments}
\subsection{Dataset and Protocol}
We evaluate on \textbf{MVTec AD}~\cite{bergmann2019mvtec}, the standard
benchmark for unsupervised industrial anomaly detection, comprising
\textbf{15} object and texture categories (5{,}354 images). Only the
\texttt{train/good} split is used for learning; all evaluation is on the
\texttt{test} split. The main metric is image-level AUROC; a global
anomaly score per image is obtained by max-pooling the six-yao
sensor responses over all spatial locations. All experiments run on CPU
(no GPU) with a frozen ResNet18 backbone at input resolution $256\times256$;
feature maps are aligned to the $/4$ grid ($64\times64$).

\subsection{Main Results}
Table~\ref{tab:main} and Fig.~\ref{fig:main} report per-class and average
image-level AUROC on all 15 categories. The pretrained six-yao sensor array (deep layer, L3)
attains an average AUROC of \textbf{0.862}, while the hand-crafted
(no-pretraining) instantiation attains \textbf{0.709}. For reference we
list the reported PaDiM (0.975) and PatchCore (0.992) averages on the same
benchmark. Our framework thus reaches a substantial fraction of
pretrained SOTA \emph{without any memory bank}, storing only a
$\sim$29\,MB template per class, and its hand-crafted variant is fully
interpretable with zero external priors.

\begin{table*}[t]
\centering
\caption{Image-level AUROC on all 15 MVTec AD categories. ``Hand-crafted''
is the no-pretraining instantiation (fixed convolutions); L1/L2/L3 are the
six-yao sensor array on frozen ResNet18 layers 1/2/3. Best per-class among
our variants in bold. PaDiM/PatchCore are reported reference numbers.}
\label{tab:main}
\small
\begin{tabular}{lcccccc}
\toprule
Class & Hand-crafted & L1 & L2 & \textbf{L3} & PaDiM & PatchCore \\
\midrule
bottle     & 0.960 & 0.991 & 1.000 & \textbf{0.998} & 0.979 & 1.000 \\
cable      & 0.628 & 0.593 & 0.772 & \textbf{0.795} & 0.886 & 0.993 \\
capsule    & 0.736 & 0.879 & \textbf{0.920} & 0.880 & 0.968 & 0.981 \\
carpet     & 0.403 & 0.797 & \textbf{0.917} & 0.860 & 0.991 & 0.997 \\
grid       & 0.719 & 0.549 & 0.598 & \textbf{0.839} & 0.945 & 0.997 \\
hazelnut   & 0.589 & 0.434 & 0.403 & \textbf{0.529} & 0.989 & 1.000 \\
leather    & 0.840 & 0.761 & \textbf{0.981} & 0.978 & 0.994 & 1.000 \\
metal\_nut & 0.728 & 0.814 & 0.783 & \textbf{0.908} & 0.952 & 1.000 \\
pill       & 0.528 & 0.726 & 0.817 & \textbf{0.836} & 0.947 & 0.988 \\
screw      & 0.535 & 0.575 & 0.617 & \textbf{0.722} & 0.752 & 0.982 \\
tile       & 0.772 & 0.697 & 0.904 & \textbf{0.932} & 0.975 & 1.000 \\
toothbrush & 0.736 & \textbf{0.989} & 0.858 & 0.839 & 0.970 & 0.997 \\
transistor & 0.815 & 0.827 & 0.802 & \textbf{0.929} & 0.970 & 1.000 \\
wood       & 0.937 & 0.842 & 0.960 & \textbf{0.973} & 0.987 & 0.993 \\
zipper     & 0.702 & 0.783 & 0.898 & \textbf{0.911} & 0.919 & 0.996 \\
\midrule
\textbf{Average} & \textbf{0.709} & 0.751 & 0.815 & \textbf{0.862} & 0.975 & 0.992 \\
\bottomrule
\end{tabular}
\end{table*}

\begin{figure*}[t]
\centering
\includegraphics[width=0.92\textwidth]{figures_paper/fig10_main_results.png}
\caption{Per-class image-level AUROC on all 15 MVTec AD categories for the
hand-crafted instantiation and the six-yao sensor array on ResNet18 layers
L1/L2/L3. Dashed lines mark the reported PaDiM (0.975) and PatchCore
(0.992) averages. Deep-layer (L3) six-yao sensors reach 0.862 on average.}
\label{fig:main}
\end{figure*}

Two findings stand out. (i) \emph{Deep layers dominate}: L3 (0.862)
clearly outperforms L2 (0.815) and L1 (0.751), consistent with the notion
that anomalies are changes in high-level structure. The shallow layer
L1 is markedly stronger only on \texttt{toothbrush} (0.989), where the
defect is a fine local lesion. (ii) \emph{Texture categories with strong
intra-class shape variation} (\texttt{hazelnut} 0.529, \texttt{cable}
0.795) and \emph{fine-grained categories} (\texttt{screw} 0.722) remain
the weakest---a genuine, interpretable limitation discussed in
Section~\ref{sec:ablation}.

\subsection{Signal-Agnosticism}
The same six-yao framework, with the extractor swapped and the criterion
fixed, spans a wide performance ladder (Table~\ref{tab:main}): hand-crafted
fixed convolution 0.709 $<$ ImageNet-pretrained ResNet 0.862. A
randomly-initialized ResNet (structure without knowledge) is
near-useless (\textbf{0.772}), confirming that the framework does not
\emph{produce} discriminative knowledge: its performance is bounded by
the information content of the incoming signal. Notably, the hand-crafted
variant is \emph{more robust} than the pretrained one on
\texttt{hazelnut} (0.589 vs.\ 0.529) and \texttt{carpet}
(0.403 vs.\ 0.860): the deep pretrained features overfit to the dominant
individual-shape variation of these categories, whereas the simple fixed
convolutions respond only to local intensity structure.

\subsection{Layer Contribution, Resolution, Online vs Frozen}
\begin{figure}[t]
\centering
\includegraphics[width=\columnwidth]{figures_paper/fig7_layer_heatmap.png}
\caption{Per-class image-level AUROC of the six-yao sensor array across
ResNet18 layers L1/L2/L3 (15 classes). Deep layers dominate on
structure-type anomalies; the shallow layer L1 helps only on fine local
defects such as \texttt{toothbrush}.}
\label{fig:layer}
\end{figure}
Layer-wise average AUROC is L1\,=\,0.751, L2\,=\,0.815, L3\,=\,0.862
(Table~\ref{tab:main}); the per-class breakdown is visualized in
Fig.~\ref{fig:layer}. A naive rank-fusion of L1\,$+$\,L3 gives 0.825,
\emph{below} L3 alone---because on several categories (\texttt{cable},
\texttt{grid}, \texttt{hazelnut}, \texttt{screw}) the shallow layer is
near-random and injects noise into the fused score. This refines an
earlier 4-class observation where L1\,$+$\,L3 fusion was optimal: with the
full benchmark, the correct design choice is to rely on the deep layer and
add a shallow layer only where it demonstrably helps (e.g., fine local
defects). Online (Welford) versus frozen normal statistics gives
indistinguishable AUROC ($<0.01$) on large classes, but online learning
enables the cold-start robustness of Section~\ref{sec:ablation}
(Fig.~\ref{fig:online}).

\begin{figure}[t]
\centering
\includegraphics[width=\columnwidth]{figures_paper/fig5_online_curve.png}
\caption{Online learning curve: image-level AUROC vs.\ number of normal
samples observed, for the six-yao sensor array. Learning is stable from
the first samples (Welford $1/n$ weighting).}
\label{fig:online}
\end{figure}

\subsection{Few-Shot Performance}
\begin{figure}[t]
\centering
\includegraphics[width=\columnwidth]{figures_paper/fig8_fewshot.png}
\caption{Few-shot performance: image-level AUROC vs.\ number of normal
training images per class for the hand-crafted and pretrained
instantiations.}
\label{fig:few}
\end{figure}
Few-shot curves for the hand-crafted and pretrained signals are shown in
Fig.~\ref{fig:few}. The pretrained instantiation is strongly data-efficient:
with only 5 normal images per class it already recovers most of its
full-data AUROC, owing to the frozen backbone carrying the bulk of the
discriminative structure. The hand-crafted instantiation degrades more
gracefully but starts lower.

%=======================================================================
\section{Ablation and Analysis}
\label{sec:ablation}
\subsection{Knowledge vs Carrier: Pretraining Ablation}
We replace the feature extractor while keeping the framework fixed:
randomly-initialized ResNet18 (structure, no knowledge), hand-crafted
fixed convolution (Sobel/Laplacian/average), hand-crafted six-yao signals
(zero pretraining), and ImageNet-pretrained ResNet18. On a common 4-class
subset (\texttt{bottle}, \texttt{tile}, \texttt{metal\_nut},
\texttt{toothbrush}) the average AUROC ladder is
\textbf{0.772} (random init) $<$ \textbf{0.799} (fixed convolution)
$<$ \textbf{0.854} (hand-crafted six-yao) $<$ \textbf{0.932}
(ImageNet-pretrained); over all 15 classes the hand-crafted-six-yao and
pretrained variants reach \textbf{0.709} and \textbf{0.862} respectively.
This demonstrates that
performance is bounded by signal information: the convolution operator is
an empty carrier, and the framework organizes rather than produces
knowledge. Notably, on categories dominated by individual-shape variation
(\texttt{hazelnut}, \texttt{carpet}) the hand-crafted signal even
\emph{outperforms} the pretrained one, because the latter's rich features
overfit the normal intra-class variance (Fig.~\ref{fig:ladder}).

\begin{figure}[t]
\centering
\includegraphics[width=\columnwidth]{figures_paper/fig4_knowledge_ladder.png}
\caption{The knowledge-vs-carrier ladder: the same
six-yao framework with different signal sources. Random-initialized ResNet
($0.772$) $<$ hand-crafted fixed convolution ($0.799$) $<$ hand-crafted
six-yao signals ($0.854$) $<$ ImageNet-pretrained ResNet ($0.932$), on a
common 4-class subset. Performance is bounded by the information content of
the signal.}
\label{fig:ladder}
\end{figure}

\subsection{Yao Variation vs Orthogonality}
An early hypothesis was that whitening the six-yao space (PCA
orthogonalization) would sharpen the anomaly signal by removing
redundancy. The controlled comparison (Table~\ref{tab:ortho}) shows the
opposite: exhaustive orthogonalization \emph{hurts} (average AUROC 0.803
$\to$ 0.731), confirming that the correlations between yao carry
discriminative information---precisely what the yao-variation view says
should be preserved, not removed. A \emph{lite} variant that retains
correlations performs on par with the raw representation (0.802). This is
a core negative result: for the six-yao representation, correlation is
\emph{information}, not redundancy.

\begin{table}[t]
\centering
\caption{Yao variation vs.\ orthogonalization (4 classes).}
\label{tab:ortho}
\begin{tabular}{lc}
\toprule
Variant & Avg.\ AUROC \\
\midrule
Raw six-yao             & 0.8033 \\
Full orthogonalization  & 0.7312 \\
Lite orthogonalization  & 0.8022 \\
Designed orthogonal     & 0.8132 \\
\bottomrule
\end{tabular}
\end{table}

\subsection{Automatic Sensitivity Learning Fails}
We tested three ways to let the framework learn \emph{which} yao matter,
using only normal data: (i) sensitive-yao weighting via the inverse
covariance $\Sigma^{-1}$, (ii) adaptive per-layer weighting, and (iii)
covariance shrinkage. All failed or were neutral. Sensitive-yao weighting
is no better than uniform weighting (e.g., \texttt{tile} 0.701 vs.\ 0.714;
\texttt{toothbrush} 0.778 vs.\ 0.767), because under pure unsupervised
learning \emph{sensitivity cannot be inferred from normal-data
statistics}; it can only come from anomaly-response signals. Covariance
shrinkage is a no-op: the estimated shrinkage coefficients are tiny
(0.01--0.05), i.e.,\ the covariances are not ill-conditioned. Adaptive
layer weighting degenerates to a near-uniform $\approx0.33$ and
underperforms the fixed L1$+$L3 choice (0.937 vs.\ 0.947).

\subsection{Threshold Rule Scan}
We scanned candidate quantile rules $q\in\{50,60,70,75,80,90\}$ for
setting the decision threshold from the model's own normal scores. No
single rule is optimal across classes: \texttt{bottle} favors a lower
quantile while \texttt{tile} favors a higher one, and the per-class
optimum varies non-monotonically with the feature signal. This is why we
adopt a single, parameter-free, robust working point: $q{=}75$. It is not
the per-class optimum, but it is a defensible compromise that requires no
per-category tuning---consistent with the zero-configuration goal.

\subsection{Cold-Start Contamination}
A key advantage of online learned statistics is robustness to a
\emph{contaminated} start. We inject $k$ defective samples
($k\in\{0,1,2,3,5\}$) at the very front of the normal stream and
measure the final AUROC after the full (clean) normal stream has been
learned. Because Welford's $1/n$ weighting dilutes early contributions,
the effect is small: on MVTec \texttt{bottle} the AUROC drops only from
0.975 to 0.969 at $k{=}5$, and on \texttt{tile} it is essentially flat
(0.878 at $k{=}0$ vs.\ 0.881 at $k{=}5$) (Table~\ref{tab:cold}).
For the small \texttt{toothbrush} class (60 normal images) the effect is
larger (0.842$\to$0.758) because the contamination represents a large
fraction of the stream---an honest, interpretable limitation
(Fig.~\ref{fig:cold}).

\begin{figure}[t]
\centering
\includegraphics[width=\columnwidth]{figures_paper/fig6_coldstart.png}
\caption{Cold-start contamination: final AUROC vs.\ the number $k$ of
defective samples injected at the front of the normal stream. Large
classes are essentially immune; the small \texttt{toothbrush} class
degrades more because the contamination dominates its short stream.}
\label{fig:cold}
\end{figure}

\begin{table}[t]
\centering
\caption{Cold-start contamination: final AUROC after injecting $k$
defective samples at the front of the normal stream.}
\label{tab:cold}
\begin{tabular}{lccccc}
\toprule
Class & $k{=}0$ & $k{=}1$ & $k{=}2$ & $k{=}3$ & $k{=}5$ \\
\midrule
bottle     & 0.975 & 0.974 & 0.974 & 0.975 & 0.969 \\
tile       & 0.878 & 0.878 & 0.884 & 0.882 & 0.881 \\
metal\_nut & 0.746 & 0.728 & 0.726 & 0.728 & 0.725 \\
toothbrush & 0.842 & 0.811 & 0.786 & 0.767 & 0.758 \\
\bottomrule
\end{tabular}
\end{table}

%=======================================================================
\section{Discussion}
\subsection{The Yao-Variation View}
Conventional statistical AD removes correlation (whitening) and measures
distance in the decorrelated space. Our results indicate the opposite is
true for the six-yao representation: correlation among yao readings is
\emph{informative}, and orthogonalizing it removes signal. This aligns with
the \emph{Yijing} view that the essence of an anomaly is \emph{change}---
the alteration of relations among the constituents---rather than mere
magnitude.

\subsection{Knowledge Versus Carrier}
The pretraining ablation isolates a conceptually important point: the
convolution operation is an \emph{empty carrier}. Randomly-initialized
ResNet features are near-useless (\textbf{0.772}), hand-crafted fixed
convolutions recover part of the signal (\textbf{0.709}), and
ImageNet-pretrained features reach \textbf{0.862}. The framework does not
\emph{produce} discriminative knowledge; it \emph{organizes} the knowledge
already present in the signal. This is why the same framework spans the
whole ladder.

\subsection{The Value of Negative Results}
We report several failures---orthogonalization, automatic sensitivity
learning via $\Sigma^{-1}$, adaptive layer weighting, and shrinkage
covariance---all of which point to the same conclusion: under pure
unsupervised learning, \emph{sensitivity cannot be inferred from
normal-data statistics}; it must come from anomaly-response signals. We
consider this a useful methodological finding rather than a mere
limitation.

\subsection{Limitations}
The pretrained instantiation is not truly prior-free: it borrows ImageNet
visual priors (though it uses no target-domain labels and no fine-tuning).
The \texttt{hazelnut} category (average AUROC 0.529) remains the weakest
for all deep-feature routes, because its 391 normal training images exhibit
large individual-shape variation: the normal template is blurred and the
defect signal is swamped. \texttt{cable} and \texttt{screw} likewise
suffer from fine, low-contrast defects that a coarse $/4$ grid may lose.
Future work includes shape/geometry sensors for shape-type defects and
higher-resolution feature grids.

\section{Conclusion}
We presented a signal-agnostic, Yijing-inspired framework for unsupervised
industrial anomaly detection, built on a six-yao sensor array and a
yao-variation criterion, with parameter-free online adaptation. The same
framework serves both hand-crafted and pretrained signals, is robust to
cold-start contamination, and is data-efficient. On the full 15-class
MVTec AD benchmark it reaches 0.862 average AUROC with pretrained features
and 0.709 with purely hand-crafted signals, using no memory bank. Future
work includes shape/geometry sensors for shape-type defects, higher-resolution
feature grids, and extension to video/sequence streams.

\section*{Funding}
This work was supported by the National Natural Science Foundation of China
(No.\ 22374086) and the National Key R\&D Program of China (No.\ 2023YFF0612100).

%=======================================================================
\bibliographystyle{IEEEtran}
\begin{thebibliography}{99}
\bibitem{bergmann2019mvtec} P. Bergmann \emph{et al.}, ``MVTec AD---A
comprehensive real-world dataset for unsupervised anomaly detection,''
in \emph{CVPR}, 2019.
\bibitem{defard2021padim} T. Defard \emph{et al.}, ``PaDiM: A patch
distribution modeling framework for anomaly detection and localization,''
in \emph{ICPR}, 2021.
\bibitem{roth2022patchcore} K. Roth \emph{et al.}, ``Towards total recall
in industrial anomaly detection,'' in \emph{CVPR}, 2022.
\bibitem{cohen2021spade} N. Cohen and Y. Hoshen, ``Sub-image anomaly
detection with deep pyramid correspondences,'' 2021.
\bibitem{he2016resnet} K. He \emph{et al.}, ``Deep residual learning for
image recognition,'' in \emph{CVPR}, 2016.
\bibitem{jeong2023winclip} J. Jeong \emph{et al.}, ``WinCLIP: Zero-/few-shot
anomaly classification and segmentation,'' in \emph{CVPR}, 2023.
\bibitem{liu2023simplenet} Z. Liu \emph{et al.}, ``SimpleNet: A simple
network for image anomaly detection and localization,'' in \emph{CVPR}, 2023.
\bibitem{zavrtanik2021draem} V. Zavrtanik \emph{et al.}, ``DRAEM---A
discriminatively trained reconstruction embedding for surface anomaly
detection,'' in \emph{ICCV}, 2021.
\bibitem{zhang2020yijing} Y. Zhang, ``The Yijing, binary coding, and Leibniz:
a computational perspective,'' \emph{Journal of Cultural Computing}, 2020.
\end{thebibliography}

\end{document}
